\section{Alignment 扩展：从 raw capability 到 useful model}

\subsection{Base model vs aligned model}

Base model 擅长完成下一 token，但并不天然适合当助手。alignment 的任务，是把 raw capability 转成真正可用、可控、可对齐人类需求的模型。

\begin{figure}[H]
\centering
\begin{tikzpicture}[node distance=1.0cm, every node/.style={font=\small, align=center}]
\node[draw, rounded corners, fill=red!8, minimum width=3.4cm, minimum height=0.8cm] (base) {base model\\raw potential};
\node[draw, rounded corners, fill=green!12, right=1.8cm of base, minimum width=3.4cm, minimum height=0.8cm] (aligned) {aligned model\\useful assistant};
\draw[-{Latex[length=3mm]}, thick] (base) -- node[above]{alignment} (aligned);
\end{tikzpicture}
\caption{alignment 的作用：把潜在能力转成可用能力}
\end{figure}

\begin{importantbox}{alignment 的三个目标}
\begin{itemize}
    \item follow instructions.
    \item control style.
    \item respect safety boundaries.
\end{itemize}
\end{importantbox}

\subsection{SFT}

SFT 的训练数据是 `(prompt, response)`。它的核心不是“再训一遍”，而是让模型明确知道：在这个 prompt 下，人类希望看到什么样的回答。

\begin{figure}[H]
\centering
\begin{tikzpicture}[node distance=1.0cm, every node/.style={font=\small, align=center}]
\node[draw, rounded corners, fill=blue!12, minimum width=3.0cm, minimum height=0.8cm] (p) {prompt};
\node[draw, rounded corners, fill=blue!20, right=1.6cm of p, minimum width=3.0cm, minimum height=0.8cm] (r) {response};
\node[draw, rounded corners, fill=blue!28, right=1.6cm of r, minimum width=3.0cm, minimum height=0.8cm] (obj) {maximize p(response | prompt)};
\draw[-{Latex[length=3mm]}, thick] (p) -- (r);
\draw[-{Latex[length=3mm]}, thick] (r) -- (obj);
\end{tikzpicture}
\caption{SFT 的基本结构：从 prompt-response pair 中学习条件分布}
\end{figure}

\begin{knowledgebox}{为什么 SFT 重要}
\begin{itemize}
    \item base model 可能知道答案，但不知道格式。
    \item SFT 把潜在能力拉到前台。
    \item 它是 instruction following 的第一步。
\end{itemize}
\end{knowledgebox}

\subsection{Preference data 与 verifiers}

进一步的对齐不再要求每个样本都有唯一标准答案，而是让模型在多个候选中选更好的那个。这个思路可以由人类偏好，也可以由形式化 verifier 或 learned verifier 支持。

\begin{figure}[H]
\centering
\begin{tikzpicture}[node distance=0.9cm, every node/.style={font=\small, align=center}]
\node[draw, rounded corners, fill=yellow!12, minimum width=3.0cm, minimum height=0.8cm] (pp) {prompt};
\node[draw, rounded corners, fill=yellow!20, right=1.4cm of pp, minimum width=2.8cm, minimum height=0.8cm] (a) {A};
\node[draw, rounded corners, fill=yellow!20, below=of a, minimum width=2.8cm, minimum height=0.8cm] (b) {B};
\node[draw, rounded corners, fill=yellow!28, right=1.8cm of a, minimum width=3.0cm, minimum height=1.1cm] (judge) {preference / verifier};
\draw[-{Latex[length=3mm]}, thick] (pp) -- (a);
\draw[-{Latex[length=3mm]}, thick] (pp) -- (b);
\draw[-{Latex[length=3mm]}, thick] (a) -- (judge);
\draw[-{Latex[length=3mm]}, thick] (b) -- (judge);
\end{tikzpicture}
\caption{偏好学习：比较多个候选回答，而不是只看单个标准答案}
\end{figure}

\begin{table}[H]
\centering
\begin{tabular}{p{3cm}p{4cm}p{5.8cm}}
\toprule
\textbf{方法} & \textbf{对象} & \textbf{直觉} \\
\midrule
PPO & policy + reward model & 把偏好当 RL 问题 \\
DPO & preference pairs & 直接从偏好对中学 policy \\
GRPO & group comparisons & 用组比较减少 value function 依赖 \\
\bottomrule
\end{tabular}
\caption{课程中会出现的三类偏好优化方法}
\end{table}

\begin{importantbox}{verifier 的价值}
对 code / math，可以有形式化 verifier；对更开放任务，可以训练 learned verifier。它们的作用都是把“更好”这件事从纯人工标注中部分解放出来。
\end{importantbox}

